% GENERATED FILE. Edit canonical lessons, then run npm run generate:books.
% canonical-source-hash: fnv1a64:8ffeb0bc5a600457
% canonical-lessons: HI-C320-samrat, HI-C320-badsah, HI-C320-simhasan, HI-C320-mukut, HI-C320-darbar

\chapter{Kings and Emperors}
\label{ch:hi-kings-and-emperors}
\hlchaptermodality{320}

\begin{chapteropening}
\textbf{By the end of this chapter:} \emph{I can tell who ruled long ago and how they lived.}

Complete the last lesson of chapter 320: I can tell who ruled long ago and how they lived.
\end{chapteropening}

\section[\texorpdfstring{\hi{सम्राट} (samrāṭ)}{सम्राट (samrāṭ)}]{\hi{सम्राट} (samrāṭ) — an emperor}

\label{lesson:HI-C320-samrat}



\begin{quote}
Before the new one: say the Hindi for a slap, then the Hindi for to twist.
\end{quote}

\subsection*{You'll want to know: \hi{सम्राट}}
\textbf{\hi{सम्राट}} — \emph{samrāṭ} — \textquotedblleft{}an emperor\textquotedblright{}.

\begin{quote}
\hi{सम्राट} \hi{के} \hi{पास} \hi{बहुत} \hi{बड़ी} \hi{सेना} \hi{थी।}
\end{quote}

\begin{quote}
\emph{samrāṭ ke pās bahut baṛī senā thī.}
\end{quote}

\begin{quote}
The emperor had a very large army.
\end{quote}

Masculine.

\begin{cousinweb}
\textbf{\hi{सम्राट}} (\emph{samrāṭ}) is Sanskrit \emph{samrāj}, \textquotedblleft{}a universal ruler\textquotedblright{} (Monier-Williams), whose nominative is \emph{samrāṭ}: a king above kings.
\end{cousinweb}

\begin{grammarlens}[title={What you've built}]
One more word for telling what happened.
\end{grammarlens}

\subsection*{Your turn}
\begin{itemize}
\raggedright
  \item \emph{Say it:} \emph{samrāṭ}
  \item \emph{Say it:} \emph{samrāṭ}, once more
  \item \emph{Say it:} \emph{samrāṭ ke pās bahut baṛī senā thī.} — The emperor had a very large army.
  \item \emph{Recall:} say the Hindi for a slap, then the Hindi for to twist, then say \emph{samrāṭ} again
\end{itemize}

\begin{tcolorbox}[breakable,colback=teal!4,colframe=teal!35!black,title={Before you move on}]
What does \hi{सम्राट} mean? (\textquotedblleft{}An emperor\textquotedblright{}.) Say it once more.
\end{tcolorbox}
\section[\texorpdfstring{\hi{बादशाह} (bādśāh)}{बादशाह (bādśāh)}]{\hi{बादशाह} (bādśāh) — a sovereign; a king of kings}

\label{lesson:HI-C320-badsah}



\begin{quote}
Before the new one: say the Hindi for to twist, then the Hindi for an emperor.
\end{quote}

\subsection*{You'll want to know: \hi{बादशाह}}
\textbf{\hi{बादशाह}} — \emph{bādśāh} — \textquotedblleft{}a sovereign; a king of kings\textquotedblright{}.

\begin{quote}
\hi{बादशाह} \hi{ने} \hi{क़िले} \hi{में} \hi{एक} \hi{बड़ी} \hi{दावत} \hi{दी।}
\end{quote}

\begin{quote}
\emph{bādśāh ne qile meṁ ek baṛī dāvat dī.}
\end{quote}

\begin{quote}
The emperor gave a great feast in the fort.
\end{quote}

Masculine.

\begin{cousinweb}
\textbf{\hi{बादशाह}} (\emph{bādśāh}) is usually traced to Persian \emph{pādšāh}, \textquotedblleft{}a king\textquotedblright{}, the title of the Mughal rulers. In stories it is the king of a Persian-style court.
\end{cousinweb}

\begin{grammarlens}[title={What you've built}]
One more word for telling what happened.
\end{grammarlens}

\subsection*{Your turn}
\begin{itemize}
\raggedright
  \item \emph{Say it:} \emph{bādśāh}
  \item \emph{Say it:} \emph{bādśāh}, once more
  \item \emph{Say it:} \emph{bādśāh ne qile meṁ ek baṛī dāvat dī.} — The emperor gave a great feast in the fort.
  \item \emph{Recall:} say the Hindi for to twist, then the Hindi for an emperor, then say \emph{bādśāh} again
\end{itemize}

\begin{tcolorbox}[breakable,colback=teal!4,colframe=teal!35!black,title={Before you move on}]
What does \hi{बादशाह} mean? (\textquotedblleft{}A sovereign; a king of kings\textquotedblright{}.) Say it once more.
\end{tcolorbox}
\section[\texorpdfstring{\hi{सिंहासन} (siṁhāsan)}{सिंहासन (siṁhāsan)}]{\hi{सिंहासन} (siṁhāsan) — a throne}

\label{lesson:HI-C320-simhasan}



\begin{quote}
Before the new one: say the Hindi for an emperor, then the Hindi for a sovereign.
\end{quote}

\subsection*{You'll want to know: \hi{सिंहासन}}
\textbf{\hi{सिंहासन}} — \emph{siṁhāsan} — \textquotedblleft{}a throne\textquotedblright{}.

\begin{quote}
\hi{राजा} \hi{सोने} \hi{के} \hi{सिंहासन} \hi{पर} \hi{बैठा।}
\end{quote}

\begin{quote}
\emph{rājā sone ke siṁhāsan par baiṭhā.}
\end{quote}

\begin{quote}
The king sat on a golden throne.
\end{quote}

Masculine.

\begin{cousinweb}
\textbf{\hi{सिंहासन}} (\emph{siṁhāsan}) is Sanskrit \emph{siṁhāsana}, \emph{siṁha}, a lion, and \emph{āsana}, a seat (Monier-Williams): the lion-seat of a king.
\end{cousinweb}

\begin{grammarlens}[title={What you've built}]
One more word for telling what happened.
\end{grammarlens}

\subsection*{Your turn}
\begin{itemize}
\raggedright
  \item \emph{Say it:} \emph{siṁhāsan}
  \item \emph{Say it:} \emph{siṁhāsan}, once more
  \item \emph{Say it:} \emph{rājā sone ke siṁhāsan par baiṭhā.} — The king sat on a golden throne.
  \item \emph{Recall:} say the Hindi for an emperor, then the Hindi for a sovereign, then say \emph{siṁhāsan} again
\end{itemize}

\begin{tcolorbox}[breakable,colback=teal!4,colframe=teal!35!black,title={Before you move on}]
What does \hi{सिंहासन} mean? (\textquotedblleft{}A throne\textquotedblright{}.) Say it once more.
\end{tcolorbox}
\section[\texorpdfstring{\hi{मुकुट} (mukuṭ)}{मुकुट (mukuṭ)}]{\hi{मुकुट} (mukuṭ) — a crown}

\label{lesson:HI-C320-mukut}



\begin{quote}
Before the new one: say the Hindi for a sovereign, then the Hindi for a throne.
\end{quote}

\subsection*{You'll want to know: \hi{मुकुट}}
\textbf{\hi{मुकुट}} — \emph{mukuṭ} — \textquotedblleft{}a crown\textquotedblright{}.

\begin{quote}
\hi{रानी} \hi{ने} \hi{मुकुट} \hi{पहना।}
\end{quote}

\begin{quote}
\emph{rānī ne mukuṭ pahnā.}
\end{quote}

\begin{quote}
The queen put on her crown.
\end{quote}

Masculine.

\begin{cousinweb}
\textbf{\hi{मुकुट}} (\emph{mukuṭ}) is Sanskrit \emph{mukuṭa}, \textquotedblleft{}a crown, a diadem\textquotedblright{} (Monier-Williams). You \textbf{\hi{पहनना}} a crown as you wear clothes.
\end{cousinweb}

\begin{grammarlens}[title={What you've built}]
One more word for telling what happened.
\end{grammarlens}

\subsection*{Your turn}
\begin{itemize}
\raggedright
  \item \emph{Say it:} \emph{mukuṭ}
  \item \emph{Say it:} \emph{mukuṭ}, once more
  \item \emph{Say it:} \emph{rānī ne mukuṭ pahnā.} — The queen put on her crown.
  \item \emph{Recall:} say the Hindi for a sovereign, then the Hindi for a throne, then say \emph{mukuṭ} again
\end{itemize}

\begin{tcolorbox}[breakable,colback=teal!4,colframe=teal!35!black,title={Before you move on}]
What does \hi{मुकुट} mean? (\textquotedblleft{}A crown\textquotedblright{}.) Say it once more.
\end{tcolorbox}
\section[\texorpdfstring{\hi{दरबार} (darbār)}{दरबार (darbār)}]{\hi{दरबार} (darbār) — a royal court}

\label{lesson:HI-C320-darbar}



\begin{quote}
Before the new one: say the Hindi for a throne, then the Hindi for a crown.
\end{quote}

\subsection*{You'll want to know: \hi{दरबार}}
\textbf{\hi{दरबार}} — \emph{darbār} — \textquotedblleft{}a royal court\textquotedblright{}.

\begin{quote}
\hi{दरबार} \hi{में} \hi{कवि} \hi{और} \hi{गायक} \hi{भी} \hi{थे।}
\end{quote}

\begin{quote}
\emph{darbār meṁ kavi aur gāyak bhī the.}
\end{quote}

\begin{quote}
There were poets and singers at the royal court too.
\end{quote}

Masculine.

\begin{cousinweb}
\textbf{\hi{दरबार}} (\emph{darbār}) is usually traced to Persian \emph{darbār}, \textquotedblleft{}a royal court\textquotedblright{}, whose first part is \emph{dar}, a door, as in \textbf{\hi{दरवाज़ा}}: those who stand at the king's door.
\end{cousinweb}

\begin{grammarlens}[title={What you've built}]
One more word for telling what happened.
\end{grammarlens}

\subsection*{Your turn}
\begin{itemize}
\raggedright
  \item \emph{Say it:} \emph{darbār}
  \item \emph{Say it:} \emph{darbār}, once more
  \item \emph{Say it:} \emph{darbār meṁ kavi aur gāyak bhī the.} — There were poets and singers at the royal court too.
  \item \emph{Recall:} say the Hindi for a throne, then the Hindi for a crown, then say \emph{darbār} again
\end{itemize}

\begin{tcolorbox}[breakable,colback=teal!4,colframe=teal!35!black,title={Before you move on}]
What does \hi{दरबार} mean? (\textquotedblleft{}A royal court\textquotedblright{}.) Say it once more.
\end{tcolorbox}
